\chapter{NixOS 的模型}
\chapterlead{NixOS 把系统状态写成表达式，由这些表达式算出软件仓库与配置。本章解释 \texttt{/nix/store}、derivation、代与闭包四个概念，它们决定了升级、回滚与清理的方式。}

\section{起步：在 store 里找一个包}

\begin{shellsession}
$ which git
/run/current-system/sw/bin/git
$ readlink -f $(which git)
/nix/store/9b2k3lq0m2fz8x1p6c4v7wj5n3h5d0a1-git-2.56.0/bin/git
$ ls /nix/store | head -3
$ ls /nix/store | wc -l
$ du -sh /nix
\end{shellsession}

\texttt{/run/current-system/sw/bin/git} 是一个符号链接，最终指向 \texttt{/nix/store} 下某个路径。路径名由两部分组成：一串哈希，和一个可读的名字。哈希由包的构建输入算出，输入相同则哈希相同，因此同一个包在不同机器上落在同一个路径。

\begin{booktable}{\textwidth}{J{0.667} J{1.333}}
\toprule
概念 & 含义 \\
\midrule
store & \texttt{/nix/store}，只读的软件仓库 \\
derivation & 一份构建配方：输入、脚本、输出路径 \\
闭包 & 一个包运行所需的全部 store 路径 \\
代 & 一次系统重建留下的完整状态，可切换 \\
\bottomrule
\end{booktable}

\section{只读仓库带来的三个结果}

第一，安装不会覆盖旧版本。升级只会在 store 里增加新路径，旧的仍在原处，正在运行的程序不受影响。第二，卸载不删除正在使用的东西，只有显式执行垃圾回收才会清理无人引用的路径。第三，可以并行安装互不兼容的版本：两个 Python 项目依赖不同版本的库，各自的环境指向各自的路径，互不干扰。

代价是磁盘占用随使用增长。同一软件的多个版本、构建过程中的中间产物都留在 store 里，需要定期清理，方法在第十五章。

\begin{figure}[H]
\centering
\figfit{%
\begin{tikzpicture}[font=\small, node distance=8mm,
  box/.style={draw=ink!45, fill=ink!5, rounded corners=1pt, inner sep=4pt, align=center},
  lab/.style={font=\scriptsize, text=muted, align=center}]
  \node[box] (expr) {配置文件\\[-1pt]{\scriptsize configuration.nix}};
  \node[box, right=of expr] (drv) {derivation\\[-1pt]{\scriptsize 构建配方}};
  \node[box, right=of drv] (store) {store 路径\\[-1pt]{\scriptsize 哈希 + 名字}};
  \node[box, right=of store] (gen) {代\\[-1pt]{\scriptsize 系统入口}};
  \draw[-{Latex[length=1.5mm]}, teal, thick] (expr) -- (drv);
  \draw[-{Latex[length=1.5mm]}, teal, thick] (drv) -- (store);
  \draw[-{Latex[length=1.5mm]}, teal, thick] (store) -- (gen);
  \node[lab, below=1mm of expr] {用户所写};
  \node[lab, below=1mm of drv] {Nix 算出};
  \node[lab, below=1mm of store] {落盘};
  \node[lab, below=1mm of gen] {可切换};
\end{tikzpicture}%
}
\caption{从配置到可用系统。每一步的产物都可复现}
\end{figure}

\section{代与切换}

每次 \texttt{nixos-rebuild switch} 生成一份新的系统配置，记为一个代。旧的代保留，因此可以回滚。

\begin{shellsession}
$ sudo nixos-rebuild list-generations
$ ls -l /nix/var/nix/profiles/system-*-link | tail -3
$ sudo /nix/var/nix/profiles/system-2-link/bin/switch-to-configuration switch
\end{shellsession}

引导菜单里也保留旧代，系统起不来时可以在启动时选择上一个。这条路径是 NixOS 在服务器上可靠的原因之一：升级失败不用进救援介质。

\section{闭包}

一个包的闭包是它运行所需的全部依赖。查看闭包可以解释两件事：为什么装一个小程序会下载很多东西，以及为什么磁盘占用比预期高。

\begin{shellsession}
$ nix path-info -Sh $(which git)
$ nix path-info -rSh $(which git) | sort -k3 -h | tail -5
$ nix why-depends $(which git) /nix/store/...-glibc-2.40
\end{shellsession}

\texttt{-r} 递归列出闭包，\texttt{-S} 显示实际占用。\texttt{nix why-depends} 回答“这个包为什么依赖那个包”，输出一条从前者到后者的路径。

\section{不可变与可变的分界}

store 只读，但系统并非全部只读。用户数据、日志、数据库、容器卷都在 store 之外，由各自的服务管理。分界线大致如下。

\begin{booktable}{\textwidth}{J{0.923} J{1.077}}
\toprule
可变 & 位置 \\
\midrule
用户文件 & \texttt{/home} \\
服务状态 & \texttt{/var/lib} \\
日志 & \texttt{/var/log/journal} \\
运行时 & \texttt{/run} \\
\bottomrule
\end{booktable}

理解这条分界，就不会试图去改 store 里的文件：改了会在下次构建时被覆盖，而且哈希不再匹配。要改变程序行为，改配置或改调用方式。

\section{练习}

\begin{enumerate}
  \item 用 \texttt{readlink -f} 找出 \texttt{bash}、\texttt{ls}、\texttt{ssh} 三个命令的 store 路径，观察命名规律。
  \item 用 \texttt{nix path-info -Sh} 查看 \texttt{git} 的闭包大小，再用 \texttt{-rSh} 找出最大的三个依赖。
  \item 用 \texttt{nix why-depends} 解释 \texttt{git} 为什么依赖 \texttt{glibc}。
  \item 查看当前的代数与创建时间，指出最近三次重建分别改了什么。
  \item 在引导菜单里找到旧代，指出它与 \texttt{list-generations} 输出中的哪一项对应。
\end{enumerate}

\begin{checkpoint}
本章结束时，应能：
\begin{itemize}
  \item 解释 store 路径的构成，说明哈希的作用；
  \item 用 store、derivation、闭包、代四个概念描述一次重建；
  \item 说明只读仓库带来的三个结果与一项代价；
  \item 用 \texttt{list-generations} 与引导菜单回滚；
  \item 分清哪些数据在 store 之外，为什么不能直接改 store 里的文件。
\end{itemize}
\end{checkpoint}
